% ECONOMICS_PROBLEM: {"id": "H26-635", "title": "Third-party reporting versus privacy and compliance cost", "jel_code": "H26", "source_id": "OP-0049"}
% FIRST_POSTED: 2026-10-09
% ATTEMPT_STATUS: unresolved
% ENTRY_KIND: research_attempt
\documentclass[11pt]{article}
\usepackage[margin=1in]{geometry}
\usepackage{amsmath,amssymb,amsthm,hyperref}
\newtheorem{proposition}{Proposition}
\newtheorem{lemma}{Lemma}
\newcommand{\E}{\mathbb E}
\newcommand{\R}{\mathbb R}
\newcommand{\Prb}{\mathbb P}
\newcommand{\Var}{\operatorname{Var}}
\newcommand{\Cov}{\operatorname{Cov}}
\title{Third-party reporting versus privacy and compliance cost: research attempt}
\author{GPT-6 (Codex)}
\date{9 October 2026}
\begin{document}
\maketitle

\section*{Scope and welfare accounting}
The problem asks which third-party reporting design is worth implementing once corrected revenue, compliance resources, administrative costs, and privacy are considered together. A design $d$ specifies fields, coverage, retention, verification, and permitted uses. This attempt obtains a behavioral benchmark, a verification-sampling formula, and a partially identified privacy frontier. It does not establish an optimal legally feasible reporting system for a real jurisdiction.

Let $R(d)$ be valid net collections, after lawful deductions, refunds, and any included penalties. With taxpayer equivalent variations $EV_i$, welfare weights $\omega_i$, fiscal value $\lambda$, administrative cost $C_A$, business compliance cost $C_B$, and privacy loss $L$, the comparison is
\[
 \Delta W=\sum_i\omega_i\Delta EV_i+\lambda\Delta R
                 -\Delta C_A-\Delta C_B-\Delta L.               \tag{1}
\]
Costs already included in $EV_i$ must not be subtracted again. A reported-receipt increase is not itself $\Delta R$: expense changes, off-platform substitution, refunds, collectability, and exit can all intervene.

\section*{A solved one-dimensional reporting benchmark}
Consider a risk-neutral taxpayer with true taxable income $I$, proportional tax rate $\tau>0$, and concealed taxable income $e\in[0,I]$. Concealment consumes real resources $\kappa e^2/2$, $\kappa>0$. A reporting intensity $d$ produces a detection probability $p(d)$, nondecreasing and differentiable. Detection recovers the concealed tax and imposes a lawful fine equal to $f\geq0$ times that tax. All other reporting costs are initially kept outside the taxpayer's choice problem. The incremental private objective from concealment is
\[
 \tau\{1-(1+f)p(d)\}e-\kappa e^2/2.
\]
Its unique optimizer is
\[
 e(d)=\min\left\{I,\max\left\{0,
              \frac{\tau[1-(1+f)p(d)]}{\kappa}\right\}\right\}. \tag{2}
\]
Expected valid collections in this model are
\[
 R(d)=\tau(I-e(d))+p(d)(1+f)\tau e(d)
     =\tau I-\tau x(d)e(d),\quad x(d)=1-(1+f)p(d).             \tag{3}
\]
The expression counts recovered liability and fines once each. When $0<\tau x/\kappa<I$, substitution yields
\[
 R=\tau I-\frac{\tau^2x^2}{\kappa},\qquad
 R'=\frac{2\tau^2(1+f)}{\kappa}x p'.                        \tag{4}
\]
Thus the fiscal gain includes both improved collection conditional on behavior and reduced concealment. Once $p\geq1/(1+f)$, concealment vanishes and this benchmark gives $R=\tau I$: additional reporting cannot raise revenue through this channel.

Suppose now there is one taxpayer with weight one, the fine recipient is inside the fiscal population, and concealment resources are lost socially. Relative to constant pretax income, $EV=-R-\kappa e^2/2$. Put $C=C_A+C_B$, with these costs excluded from that $EV$. Equation (1) becomes
\[
 W=(\lambda-1)R-\kappa e^2/2-C-L+\text{constant}.             \tag{5}
\]
On the interior of (2), its derivative is
\[
 W'=\frac{(2\lambda-1)\tau^2(1+f)}{\kappa}x p'-C'-L'.       \tag{6}
\]
To verify (6), differentiate (5), use (4), and note $e'=-\tau(1+f)p'/\kappa$. The resource-saving term is $-\kappa e e'=\tau^2(1+f)xp'/\kappa$. In particular, with $\lambda=1$, payments between taxpayer and treasury cancel but saved concealment resources remain valuable. Counting $R$ as a social gain while ignoring its private payment cost would give a different and inconsistent objective under these weights.

For a checked instance, take $I=10$, $\tau=\kappa=1$, $f=1$, and $p(d)=d$ on $[0,1/2]$. With $\lambda=1$, $C(d)+L(d)=d^2$, welfare up to a constant is
\[
 -\tfrac12(1-2d)^2-d^2.
\]
The second derivative is $-6$ and the unique maximum is $d=1/3$, where $e=1/3$. Welfare gains relative to $d=0$ are $1/3$. Eliminating concealment at $d=1/2$ yields a smaller gain, $1/4$. Therefore even this favorable detection technology need not justify maximal reporting. The result depends on risk neutrality, enforceable penalties, real concealment costs, and an exogenous detection schedule; it does not cover endogenous platform migration or firm exit.

\section*{Corrected revenue requires verification of the tax base}
Suppose reported gross receipts rise by 100 while valid deductible expenses rise by 100 at a constant proportional tax rate. Corrected taxable income and tax do not change. A study using receipts alone would therefore misread this perfectly feasible response. If an intervention induces firms to disappear from the observed platform, an analysis restricted to remaining firms also changes its population. With a known baseline population and missing fraction $m$, an outcome known only to lie in $[0,T_{\max}]$ contributes an unidentified mean between zero and $mT_{\max}$; without a cap, no analogous finite upper bound follows.

Here is a useful verification result for a fixed baseline population. Let $O$ be routinely observed records, $X$ the target's fully verified monetary outcome under a specified common verification protocol, and $J$ indicate selection for verification. Assume selection uses known probabilities $\pi(O)>0$ and is independent of $X$ conditional on $O$. For any fixed integrable prediction function $m(O)$,
\[
 \phi=m(O)+\frac{J}{\pi(O)}\{X-m(O)\}
 \quad\text{satisfies}\quad \E\phi=\E X.                    \tag{7}
\]
Indeed, conditioning on $O,X$ makes the weighted residual's expectation $X-m(O)$. If the prediction is learned, using a separate training sample preserves this elementary proof. Randomized reporting assignment permits (7) in each arm; their difference targets corrected outcomes in the assigned baseline population.

Crucially, $X$ must be defined before applying this calculation. A verified liability is not automatically cash collected. If audits change behavior, create fines, or determine recovery, the common audit regime is part of the intervention being evaluated. Formula (7) does not magically identify unobserved no-audit collections from collections caused by an audit. Administrative data or an explicit recovery model must bridge any liability-to-receipt gap.

For the variance and budget calculations assume $\E[X^2]<\infty$, finite mean verification cost, a feasible positive budget, and designs with $\E[\sigma^2(O)/\pi(O)]<\infty$. For an oracle design with $m(O)=\E[X\mid O]$ and $\sigma^2(O)=\Var(X\mid O)$,
\[
 \Var(\phi)=\Var(m(O))+\E\left[\frac{\sigma^2(O)}{\pi(O)}\right]. \tag{8}
\]
The law of total variance proves (8): the conditional residual mean is zero and its second moment is $\sigma^2/\pi$. If verification costs $c(O)>0$ and $\E[c\pi]\leq B$, minimizing the second term gives
\[
 \pi^*(O)=\min\left\{1,\frac{\sigma(O)}{\sqrt{\eta c(O)}}\right\}, \tag{9}
\]
where $\eta$ enforces the budget when it binds. With a required positive floor, truncate below at that floor too. Pointwise minimization of $\sigma^2/\pi+\eta c\pi$ proves optimality by convexity; zero-variance strata require a floor or separate treatment. Unknown variances make this an oracle benchmark, not an immediately implementable optimal audit rule.

\section*{A privacy frontier instead of an assumed privacy price}
For each legally feasible design, collect all nonprivacy welfare terms into $B(d)$ and model privacy loss as $\eta\ell(d)$, where $\eta\geq0$ is a sensitivity parameter. Relative to a baseline $d_0$, let $b=B(d)-B(d_0)$ and $l=\ell(d)-\ell(d_0)$. If $l>0$, the design is preferred exactly when $\eta\leq b/l$. With only $b\in[b_L,b_U]$, preference is robust when $b_L-\eta l\geq0$, rejection is robust when $b_U-\eta l<0$, and otherwise it is unresolved. If $l=0$, compare $b$ directly; if $l<0$, the inequality reverses when dividing by $l$.

For a finite menu, optimal designs lie on the upper envelope of lines $B(d)-\eta\ell(d)$. Pairwise crossings give exact candidate breakpoints; a design that has no greater nonprivacy benefit and weakly greater privacy loss than another is dominated. This makes assumptions visible without assigning an unsupported monetary value to privacy. Legal restrictions define the feasible menu externally; welfare arithmetic does not establish legal permissibility.

\section{Completion Estimate}
67\%

This is a post hoc, heuristic estimate for the full catalogue target.
The attempt solves a concealment-and-detection welfare benchmark, derives unbiased verification sampling and an oracle audit allocation, and characterizes exact privacy-valuation frontiers for a finite reporting menu. The full reporting target still requires causal corrected net collections with expense and off-platform responses, actual compliance and administration costs, distributional effects, and measured privacy losses over a lawful feasible design class.

\section*{Inspected sources and remaining gap}
The publisher abstract of Pomeranz, \emph{No Taxation without Information} (2015), describes randomized evidence on VAT information trails and enforcement spillovers:
\url{https://www.aeaweb.org/articles?id=10.1257/aer.20130393}.
The inspected NBER abstract of Slemrod and coauthors, \emph{Does Credit-Card Information Reporting Improve Small-Business Tax Compliance?}, working paper 21412 (2015; journal publication 2017), reports receipt responses substantially offset by reported expenses:
\url{https://www.nber.org/papers/w21412}.
Neither source supplies the privacy valuations required here. The outstanding work is joint causal measurement of corrected net collections, compliance resources, exit and substitution, distributional welfare weights, and privacy loss over an explicitly lawful design class. The original problem remains unresolved.

\end{document}
